% CERT rel=Theta f=sqrt(n) g=n**(1/2) c1=1 c2=1 n0=1
\ej{8}{Método maestro para $T(n)=2T(n/4)+\sqrt{n}$.}

Suponemos $n=4^h$, con $h\in\N\cup\{0\}$, para que las divisiones lleguen exactamente a $1$, y un costo base $T(1)>0$ constante.

\begin{enumerate}
\item \textbf{Identificación.} La forma del Teorema Maestro es $T(n)=aT(n/b)+f(n)$. En esta recurrencia:
\[
a=2,\qquad b=4,\qquad f(n)=\sqrt{n}.
\]
Se cumplen $a\ge1$ y $b>1$.

\item \textbf{Comparación.} Como $4^{1/2}=2$, por definición de logaritmo $\log_4 2=1/2$. Por tanto,
\[
n^{\log_b a}=n^{\log_4 2}=n^{1/2}=\sqrt{n}=f(n).
\]
Esta igualdad prueba $f(n)\in\Th{n^{\log_b a}}$: en la definición de $\Theta$ podemos tomar $c_1=c_2=1$ y $n_0=1$.

\item \textbf{Aplicación del caso 2.} El caso 2 del Teorema Maestro establece que, si $a\ge1$, $b>1$ y $f(n)\in\Th{n^{\log_b a}}$, entonces
\[
T(n)\in\Th{n^{\log_b a}\log n}.
\]
Las tres hipótesis quedaron verificadas. Este caso no requiere $\varepsilon>0$ ni condición de regularidad. Tomando $\log=\log_2$ y sustituyendo $n^{\log_b a}=\sqrt{n}$, obtenemos
\[
T(n)\in\Th{\sqrt{n}\log_2 n}.
\]
\end{enumerate}

Concluimos que $\boxed{T(n)\in\Th{\sqrt{n}\log n}}$, para $n$ potencia de $4$.
\fin
